\subsection*{全体概要：この章で学ぶこと}
\addcontentsline{toc}{subsection}{全体概要：この章で学ぶこと}

この章は、ソフトウェア開発を「技術だけの活動」と見なさず、「限られた資源を使って価値を作る活動」として扱う。開発するかしないか、買うか作るか、保守するか置き換えるか、品質をどこまで高めるか、どの要求を優先するかといった判断には、必ず経済的な側面がある。

到達目標 この章を読み終えると、次のことができるようになる。
\begin{itemize}
  \item ソフトウェア工学経済が、金銭だけでなく価値・リスク・無形資産を扱うことを説明できる。
  \item 提案、キャッシュフロー、時間価値、等価性、比較基準、代替案を区別できる。
  \item 工学的意思決定の流れを、問題理解、候補案、選択基準、評価、選択、監視として説明できる。
  \item 営利、非営利、現在時点、複数属性の意思決定技法を使い分けられる。
  \item 見積りが意思決定のための予測であり、不確実性を伴うことを説明できる。
  \item 会計、費用、財務、効率、有効性、生産性、製品ライフサイクルなどの関連概念を説明できる。
\end{itemize}

\par\smallskip
\begin{center}
\centering
\IfFileExists{assets/generated_figures/ch15/fig-ch15-01.png}{\includegraphics[width=0.96\linewidth]{assets/generated_figures/ch15/fig-ch15-01.png}}{\fbox{\parbox[c][48mm][c]{0.9\linewidth}{\centering 画像生成待ち\\15 章の全体地図}}}
\par\smallskip
\refstepcounter{figure}\label{fig:ch15-01}図 \thefigure: 15 章の全体地図
\end{center}
図 \thefigure は、15 章の全体地図を視覚的に整理したものである。
\par\smallskip
